\documentclass[a4paper,10pt]{article}
\usepackage[utf8]{inputenc}
\usepackage{graphicx}
\usepackage{subfigure}
\usepackage{fancyhdr}
\usepackage{listings}
\usepackage{ gensymb }


%opening

\pagestyle{fancy} %eigener Seitenstil
\fancyhf{} %alle Kopf- und Fußzeilenfelder bereinigen
\fancyhead[L]{SEP Abgabegespräch SS2015/16} %Kopfzeile links
\fancyhead[C]{} %zentrierte Kopfzeile
\fancyhead[R]{\includegraphics[width=1cm]{tugraz_bw}} %Kopfzeile rechts
\renewcommand{\headrulewidth}{0.4pt} %obere Trennlinie
\fancyfoot[C]{} %Seitennummer
\renewcommand{\footrulewidth}{0.4pt} %untere Trennlinie


\lstset{language=C++, basicstyle=\ttfamily\tiny}  

\begin{document}

\section*{Spaceship - database}

Es wird ein kleines Raumschiff - Verwaltungstool entwickelt.
Es gibt 2 Arten von Schiffen: Schlachtschiffe und Frachtschiffe. 
Jedes Raumschiff besitzt einen Namen, Hitpoints und eine Geschwindigkeit.
Zusätzlich besitzen Schlachtschiffe eine gewisse Feuerkraft, während Frachtschiffe eine Güterkapazität haben.

Das Programm kennt 3 Befehle: add, list und quit.
Add fügt ein Raumschiff zu einer internen Liste hinzu, während list alle bisher eingegebenen Raumschiffe ausgibt. Quit beendet das Programm.

\subsection*{Add}
\begin{verbatim}
> add [type] [name] [hitpoints] [speed] [capcity/weapons]

> add battelship saratoga 75.5 52.3 210.5
> add freithship pasteur 250.3 32.4 100.5
\end{verbatim}



\subsubsection*{Ausgabe}
\begin{verbatim}
> add battleship saratoga 100.5 55.3 200
> add freightship pasteur 55.4 99.3 200
> list
Battleship: saratoga
Hitpoints: 100.5
Speed: 55.3
Weapons Power: 200

Freightship: pasteur
Hitpoints: 55.4
Speed: 99.3
Capacity: 200

> quit
\end{verbatim}

\subsection*{Aufruf}
\texttt{./ag}

\subsection*{Bewertung}
\begin{itemize}
 \item Dokumentation und Programmierstil: 1 Punkt
 \item Einhalten der OOP Konzepte: 1 Punkt
 \item Struktur \& Korrektheit: 1 Punkt
\end{itemize}

\newpage


\end{document}
